\documentclass{article}

\usepackage{amsmath}
\usepackage{hyperref}
\usepackage{amsthm}
\usepackage{amssymb}
\usepackage{bbm}
\usepackage{fancyhdr}
\usepackage{listings}
\usepackage{cite}
\usepackage{graphicx}
\usepackage{enumitem}
\usepackage[margin=1cm]{caption}
\usepackage{subcaption}
\usepackage{tcolorbox}
\usepackage{color}
\definecolor{editorGray}{rgb}{0.95, 0.95, 0.95}


\lstset{%
    % Basic design
    backgroundcolor=\color{editorGray},
    basicstyle={\small\ttfamily},   
    frame=l,
    % Line numbers
    xleftmargin={0.75cm},
    numbers=left,
    stepnumber=1,
    firstnumber=1,
    numberfirstline=false,
    }
\lstset{
    literate={~} {$\sim$}{1}
}

\newenvironment{claim}[1]{\par\noindent\underline{Claim:}\space#1}{}
\newenvironment{claimproof}[1]{\par\noindent\underline{Proof:}\space#1}{\hfill $\blacksquare$}

\oddsidemargin 0in \evensidemargin 0in
\topmargin -0.5in \headheight 0.25in \headsep 0.25in
\textwidth 6.5in \textheight 9in
\parskip 6pt \parindent 0in \footskip 20pt

% set the header up
\fancyhead{}
\fancyhead[L]{Stanford Aeronautics \& Astronautics}
\fancyhead[R]{Fall 2019}

%%%%%%%%%%%%%%%%%%%%%%%%%%
\renewcommand\headrulewidth{0.4pt}
\setlength\headheight{15pt}
\usepackage{color}
\usepackage{amsmath}
\usepackage{amssymb}
\usepackage{graphicx}
\usepackage{comment,xspace}
\usepackage{hyperref}
%\usepackage{algorithm,algorithmic} 
\usepackage{fancybox}

%\newcommand{\todo}[1]{\par\noindent{\color{red}\raggedright\sc{#1}
%    \par\marginpar{\Large \bf $\star$}}}

%%%%%
%% If you use a font encoding package, please enter it here, i.e.,
%  \usepackage{T1enc}

%% How many levels of section head would you like numbered?
%% 0= no section numbers, 1= section, 2= subsection, 3= subsubsection
%%==>>
\setcounter{secnumdepth}{3}


%%For margin comments
%\newcommand{\todomar}[1]{\marginpar{\tiny\color{red}#1}}
%% Math defs

%For theorems
\newtheorem{assump}{Assumption}
\newtheorem{theorem}{Theorem}[section]
\newtheorem{definition}[theorem]{Definition}
\newtheorem{problem}{Problem}
\newtheorem{lemma}[theorem]{Lemma}
\newtheorem{corollary}[theorem]{Corollary}
\newtheorem{remark}[theorem]{Remark}

\newcommand{\real}{{\mathbb{R}}}
\newcommand{\reals}{\real}
\newcommand{\sphere}{{\mathbb{S}}}
\renewcommand{\natural}{{\mathbb{N}}}
\newcommand{\naturals}{\natural}

\newcommand{\poiss}{\text{Poisson}}
\newcommand{\xfree}{\mathcal X_{\text{free}}}
\newcommand{\xobs}{\mathcal X_{\text{obs}}}
\newcommand{\xgoal}{\mathcal X_{\text{goal}}}
\newcommand{\FMT}{$\text{FMT}^*\, $}


\newcommand{\eps}{\varepsilon}
\newcommand{\He}{H_{\eps}}
\newcommand{\co}{\operatorname{co}}
\newcommand{\ov}{\overline}
\newcommand{\sign}{\operatorname{sign}}
\newcommand{\vers}{\operatorname{vers}}
\newcommand{\union}{\cup}

\newcommand{\map}[3]{#1: #2 \rightarrow #3}
\newcommand{\setdef}[2]{\left\{#1 \; | \; #2\right\}}
\newcommand{\until}[1]{\{1,\dots,#1\}}
\newcommand{\proj}[1]{\operatorname{proj}_{#1}}
\newcommand{\Ve}{\operatorname{Ve}}
\newcommand{\pder}[2]{\frac{\partial #1}{\partial #2}}
\newcommand{\tder}[2]{\frac{d #1}{d #2}}
%\newcommand{\margin}[1]{\marginpar{\tiny\ttfamily#1}}

\newcommand{\subscr}[2]{#1_{\textup{#2}}}

\newcommand{\TSP}{\ensuremath{\operatorname{TSP}}}
\newcommand{\MSP}{\ensuremath{\operatorname{MSP}}}
\newcommand{\MST}{\ensuremath{\operatorname{MST}}}
\newcommand{\OSS}{\ensuremath{\operatorname{OSS}}}
\newcommand{\sRH}{\textup{sRH}\xspace}
\newcommand{\mRH}{\textup{mRH}\xspace}
\newcommand{\mRHFG}{\textup{mRHFG}\xspace}

\newcommand{\card}{\operatorname{card}}

\newcommand{\Vol}{\operatorname{Vol}}
\newcommand{\Area}{\operatorname{Area}}
\newcommand{\limn}{\lim_{n \rightarrow \infty}}
\newcommand{\limni}[1]{\lim_{n_{#1} \rightarrow \infty}}
\newcommand{\diam}{\operatorname{diam}}



%\newcommand{\proof}{\noindent{\bf Proof: \indent}}
%\newcommand{\proofover}{\hfill\vrule height8pt width6pt depth
%                0pt\newline}

\newcommand{\argmin}{\operatornamewithlimits{argmin}}
\newcommand{\reg}{$^{\textrm{\tiny\textregistered}}$}
\newcommand{\liml}{\lim_{\lambda \rightarrow +\infty}}



\newcommand{\p}[1]{\mbox{$\mathbb{P}\left(#1\right)$}} 
\newcommand{\expectation}[1]{\mbox{$\mathbb{E}\left[#1\right]$}} 
\newcommand{\probcond}[2]{\mbox{$\mathbb{P}\left(#1 \,| \, #2\right)$}} 
\newcommand{\condexpectation}[2]{\mbox{$\mathbb{E}\left(#1| #2\right)$}} 
\newcommand{\iteratedexpectation}[2]{\mbox{$\mathbb{E}\left(\mathbb{E}\left(#1| #2\right)\right)$}} 

% Added for KFMT
\newcommand{\LL}{\mathcal{L}}
\newcommand{\M}{\mathcal{M}}
\newcommand{\Mobs}{\mathcal{M}_{\text{obs}}}
\newcommand{\Mfree}{\mathcal{M}_{\text{free}}}
\newcommand{\Mgoal}{\mathcal{M}_{\text{goal}}}
\newcommand{\tx}{\tilde x}
\newcommand{\cl}{\operatorname{cl}}
\newcommand{\VF}{\operatorname{VF}}
\newcommand{\SF}{\texttt{SampleFree}}
\newcommand{\R}{\mathbb{R}}
\newcommand{\N}{\mathbb{N}}
\newcommand{\emin}{\epsilon_{\min}}
\newcommand{\amin}{a_{\min}}
\newcommand{\Amax}{A_{\max}}
\newcommand{\xinit}{x_{\text{init}}}
\newcommand{\pol}{\text{pol}}
\newcommand{\inte}{\text{int}}
\newcommand{\boxw}{\text{Box}^w}
\newcommand{\nti}{n\to\infty}
\newcommand{\ti}{\to\infty}
\newcommand{\X}{\mathcal{X}}
\newcommand{\KFMT}{$\text{KFMT}^*\, $}

% AA 274 specific
\newcommand{\writtenQ}{\raisebox{-1ex}{\includegraphics[scale=0.6]{write.png}}~}
\newcommand{\codeQ}{\raisebox{-1ex}{\includegraphics[scale=0.6]{code.png}}~}

\usepackage{xparse}
\NewDocumentCommand{\codeword}{v}{%
\texttt{#1}%
}

\usepackage{xcolor}
\definecolor{friendlybg}{HTML}{f0f0f0}
\usepackage{minted}
\setminted{
    style=friendly,
    bgcolor=friendlybg
}


\usepackage{outlines}

\usepackage{xparse}
\NewDocumentCommand{\codeword}{v}{%
\texttt{\textcolor{blue}{#1}}%
}
\usepackage{gensymb}

\newcommand{\ssmargin}[2]{{\color{blue}#1}{\marginpar{\color{blue}\raggedright\scriptsize [SS] #2 \par}}}


\setlength{\parindent}{0in}

\title{AA 274: Principles of Robotic Autonomy \\ Section 1: VM Installation, Git, Python}
\date{}

\begin{document}

\maketitle
\pagestyle{fancy}

Our goals for this section: \begin{enumerate}
	\item Install Ubuntu virtual machine locally.
	\item Learn how to use Git to pull assignments.
	\item Start working with Python and Numpy.
\end{enumerate}


\section{Virtual Machine Installation}
For this class, we recommend using a virtual machine (VM) that we provide which has all the
required software preinstalled.
Alternatively, you can natively install Linux alongside your normal OS as a dual boot and install the required
packages yourself. This may give better performance, especially for graphics rendering. However, as there
can be driver issues and other edge cases that can come up when installing linux on hardware, we can only
officially support the VM in terms of providing help with your setup.
To set up the VM on your computer, follow the instructions here: 

\url{https://github.com/PrinciplesofRobotAutonomy/aa274-docker}

\section{Using Git}
Git is a source control tool that allows us to share code with you. To obtain code for this section, type the following into your terminal:

 \begin{lstlisting}[escapechar=|]
 git clone https://github.com/StanfordASL/AA274_SECTION
\end{lstlisting}

You will use this folder to store all section material for the remainder of the class. So, to update this repo at the start of each section, type: 

 \begin{lstlisting}
 git pull
\end{lstlisting}

\section{Python}
Since our class is composed of students from AA, EE/ME, and CS, we will not assume that you have comprehensive Python knowledge. Therefore, the main purpose of this section is to get your Python coding skills spun up so you can work on the homework. If you know this material already, please help someone who doesn't know it as well.

In order to complete this part of our section, please switch to the scripts included in the code for this section. Remember, in this course we are using Python 2.7. If you do not have Python 2.7 installed on your computer, or would rather not disrupt your existing Python workspace, feel free to use the VM you just installed! You can start up a terminal with

 \begin{lstlisting}
./run.sh bash
\end{lstlisting}

Alternatively, you can directly launch a Python 2.7 REPL (containing all of the necessary packages) using the following

 \begin{lstlisting}
./run.sh python
\end{lstlisting}

Once you have had a look through the code, played around with Python a bit, and are aware of the capabilities that Python and its packages offer, complete the following problems:
{\bf
\begin{enumerate}
\item Define a $\sin$ function using NumPy
\item Find the minimum of the function using SciPy
\item Integrate the function from $[0, 1]$ using SciPy
\item Plot the function using Matplotlib from $[0, 2\pi]$
\end{enumerate}
}
Once you have done this, please submit your results \textit{and code} in one writeup file on Gradescope.

\end{document}
